\NSPage{115}%
\begin{EESegment}{EE-TGT-P115-001}
For the fixed background and initialization, all derivatives through order
\NSi{NS-F-P115-001}{m} are bounded by
\NSi{NS-F-P115-002}{C_mq^{-K_m}(1+|\log q|)^{P_m}}. Finally, \eqref{eq:9.18} is exactly
\eqref{eq:5.33} for the fixed differential polynomial
\end{EESegment}
\NSFormula{NS-F-P115-003}{display}{none}
\[
\mathcal F(u,p)=\mathscr R(u,p)=\partial_tu+(u\mathbin{\cdot}\nabla)u-\Delta u+\nabla p,
\]
\begin{EESegment}{EE-TGT-P115-002}
with \NSi{NS-F-P115-004}{\rho_j=h\sigma_j\to\infty}. The estimate \eqref{eq:9.18} concerns
the actual residual at each point after the auxiliary variables have been set equal to
\NSi{NS-F-P115-005}{Y(r,t)}. At any finite stage, the flat term is bounded uniformly over the
infinitely many dyadic labels. Cutting off the finite initialization adds errors only away from
\NSi{NS-F-P115-006}{q=0}, and those errors satisfy the same fixed-stage flatness condition. The
residual estimates hold uniformly over the whole local domain. To see this, choose one fixed,
bounded interval in \NSi{NS-F-P115-007}{X} that contains every slow support and every annular
correction support. The estimates above and \eqref{eq:5.41} apply on that interval. Outside it,
every finite state is exactly the exterior heat field from \eqref{eq:4.29}, together with the
centrifugal pressure, so its residual is zero there.
\end{EESegment}

\begin{EESegment}{EE-TGT-P115-003}
\noindent\textbf{Step 2: Sum the fields and check their Cartesian regularity.}
Lemma~\ref{lem:5.4} now gives potentials whose locally finite sum defines a velocity with
\NSi{NS-F-P115-008}{\operatorname{div}u_{\mathrm{loc}}=0}, and
\end{EESegment}
\NSFormula{NS-F-P115-009}{display}{9.20}
\begin{equation}
\forall m,N\qquad |\mathscr R(u_{\mathrm{loc}},p_{\mathrm{loc}})|_m=O(q^N)
\qquad(q\downarrow0).
\tag{9.20}\label{eq:9.20}
\end{equation}
\begin{EESegment}{EE-TGT-P115-004}
The tail estimate compares the field just built with one fixed finite state and that state's flat
remainder. The flat residuals separated at the finite stages are controlled through this comparison,
not added together. To identify the potentials and pressure of the summed fields, let
\NSi{NS-F-P115-010}{\mathcal A_0,\mathcal B_0e_\theta,p_0} be the representatives of
\NSi{NS-F-P115-011}{U_0}. These representatives are the realized slow base from
\eqref{eq:5.27}, together with the finite initialization and its pressure. Define
\end{EESegment}
\NSFormula{NS-F-P115-012}{display}{9.21}
\begin{equation}
\begin{aligned}
\mathcal A&=\mathcal A_0+\sum_{j\geq1}\chi(a_jq)A_j,\\
\mathcal Be_\theta&=\mathcal B_0e_\theta+\sum_{j\geq1}\chi(a_jq)B_j,\\
p_{\mathrm{loc}}&=p_0+\sum_{j\geq1}\chi(a_jq)p_j,
&u_{\mathrm{loc}}&=\operatorname{curl}\mathcal A+\mathcal Be_\theta.
\end{aligned}
\tag{9.21}\label{eq:9.21}
\end{equation}
\begin{EESegment}{EE-TGT-P115-005}
These formulas give the required vector potential, azimuthal field, and pressure. At the axis, the
base Stokes streamfunctions are \NSi{NS-F-P115-013}{r^2} times a smooth function of
\NSi{NS-F-P115-014}{(r^2,z,t)}, so the base vector potential is smooth in Cartesian coordinates.
The base angular field is \NSi{NS-F-P115-015}{r} times such a scalar, times
\NSi{NS-F-P115-016}{e_\theta}. For each \NSi{NS-F-P115-017}{t<1}, every annular representative
vanishes on a neighborhood of the axis. Equation~\eqref{eq:9.21} therefore proves
Theorem~\ref{thm:3.1}\textup{(i)}.
\end{EESegment}

\begin{EESegment}{EE-TGT-P115-006}
\noindent\textbf{Step 3: Check one-sided regularity away from the singular point.} The endpoint
regularity still has to be checked away from \NSi{NS-F-P115-018}{q=0}. Fix a compact spatial set
and numbers \NSi{NS-F-P115-019}{0<c<c'<q_*}. Only finitely many terms in the expansion in powers
of \NSi{NS-F-P115-020}{q^{2h}} have nonzero cutoff factors on a neighborhood of
\NSi{NS-F-P115-021}{c\leq q\leq c'}, and only finitely many correction stages
do, because both cutoff sequences tend to infinity. On the chosen compact set there are also only
finitely many dyadic bands and slow labels. By Theorem~\ref{thm:4.6} and
Proposition~\ref{prop:5.5}, the base profiles and their Stokes streamfunctions have bounds for
derivatives of every fixed order on the closed parameter range
\NSi{NS-F-P115-022}{-1\leq\eta\leq1}. The geometry chooses its labels, representatives, rounded
frequencies, and enlarged local torus rectangles from the closed range
\NSi{NS-F-P115-023}{\tau\geq0}, and it keeps those discrete choices fixed as
\NSi{NS-F-P115-024}{\tau\to0}. The pulse inverse follows a path on which the physical slow
variables stay fixed and \NSi{NS-F-P115-025}{\tau\geq0}. Differentiate its fixed linear ODE to
get the bounds at every derivative order. Combine these bounds with the quotient bound
\eqref{eq:7.33}, which uses the fixed positive primary amplitudes, and with the radial and temporal
mean maps. The result is a finite one-sided bound for every fixed derivative of each wave
potential, mean
\NSPage{116}%
vector potential, azimuthal mean component, and pressure whose cutoff is nonzero on this
neighborhood. Their smooth extensions by zero cover every boundary of their supports.
\end{EESegment}

\begin{EESegment}{EE-TGT-P116-001}
The denominator in the coordinate conversion is
\NSi{NS-F-P116-001}{L\geq1-2h>0}, so the conversion keeps all these fixed-order derivative bounds
on \NSi{NS-F-P116-002}{q\geq c}. On the annular supports,
\NSi{NS-F-P116-003}{r} stays away from zero. At the axis, use the Cartesian factors described
above. The finite sum then has uniformly bounded Cartesian space--time derivatives of every order
all the way to \NSi{NS-F-P116-004}{\tau=0}, including inside the active shell. Bound one more
time derivative and apply the fundamental theorem of calculus. This gives a uniform one-sided
limit for every derivative. Applying the same theorem on compact spatial sets shows that the limits
agree with spatial differentiation. Integrating in time shows that the limits of consecutive time
derivatives agree with one another. This proves Theorem~\ref{thm:3.1}\textup{(ii)} for
\NSi{NS-F-P116-005}{\mathcal A,\mathcal Be_\theta,p_{\mathrm{loc}}} themselves.
\end{EESegment}

\begin{EESegment}{EE-TGT-P116-002}
\noindent\textbf{Step 4: Check the exterior heat field and the flat residual.} Every annular wave
or mean correction potential, every azimuthal mean component, and every pressure correction
vanishes beyond \NSi{NS-F-P116-006}{X_b}. The leading profile and every coefficient in the
expansion in powers of \NSi{NS-F-P116-007}{q^{2h}} have zero total axial moment. Their Stokes
streamfunctions therefore vanish in the same exterior region. It follows that, for one fixed
\NSi{NS-F-P116-008}{X_{\mathrm{ext}}} beyond all the slow supports,
\NSi{NS-F-P116-009}{\mathcal A=0} and \NSi{NS-F-P116-010}{\mathcal B=K}. The leading-profile
pressure is normalized to vanish at radial infinity, and all positive-order pressure corrections
have compact support. The pressure there is therefore exactly
\NSi{NS-F-P116-011}{-\int_r^\infty K(\rho,\tau)^2\,d\rho/\rho}. Using \eqref{eq:4.29}, define
\end{EESegment}
\NSFormula{NS-F-P116-012}{display}{none}
\[
\mathcal H_{\mathrm{ext}}(s)=2^Ac_\infty\mathcal H(4s),
\qquad K(r,\tau)=r^{-1-2h}\mathcal H_{\mathrm{ext}}(\tau/r^2).
\]
\begin{EESegment}{EE-TGT-P116-003}
This is \eqref{eq:3.5}. For every integer \NSi{NS-F-P116-013}{m\geq0}, the integral formula
\eqref{eq:A.34} gives
\end{EESegment}
\NSFormula{NS-F-P116-014}{display}{none}
\[
\sup_{s\geq0}|\mathcal H_{\mathrm{ext}}^{(m)}(s)|
\leq2^Ac_\infty4^m(h)_m(1+h)_m<\infty,
\]
\begin{EESegment}{EE-TGT-P116-004}
where \NSi{NS-F-P116-015}{(b)_m} is the rising factorial, the product of m factors starting
at b and increasing by one each time, and
\NSi{NS-F-P116-016}{(b)_0=1}. The bound follows because the factor
\NSi{NS-F-P116-017}{(1+Zv)^{-h-m}} in that integral is at most one whenever
\NSi{NS-F-P116-018}{Z,v\geq0}. The required derivative bounds therefore hold on the whole
nonnegative half-line. Lemma~\ref{lem:A.6} says that \NSi{NS-F-P116-019}{K} solves the radial
heat equation, and its centrifugal pressure makes the Navier--Stokes residual exactly zero in the
exterior. Combine this with the estimates for the background residual on compact sets of profile
variables in \eqref{eq:5.41} and \eqref{eq:9.20}. This proves all of
Theorem~\ref{thm:3.1}\textup{(iii)}.
\end{EESegment}

\begin{EESegment}{EE-TGT-P116-005}
\noindent\textbf{Step 5: Check the growth of the inner velocity.} Finally, choose any fixed
\NSi{NS-F-P116-020}{X_{\mathrm{in}}\in(0,X_a)} at which the positive leading azimuthal velocity
is \NSi{NS-F-P116-021}{e_0=E_0(X_{\mathrm{in}},0)>0}. Every annular correction vanishes there.
The realized expansion in powers of \NSi{NS-F-P116-022}{q^{2h}} in \eqref{eq:5.1} has its
leading value plus a remainder \NSi{NS-F-P116-023}{O(q^{2h})} after the velocity factor
\NSi{NS-F-P116-024}{q^{-A}} is removed, at the fixed point
\NSi{NS-F-P116-025}{(X,\eta)=(X_{\mathrm{in}},0)} in profile coordinates. At
\NSi{NS-F-P116-026}{z=0}, the coordinates satisfy
\NSi{NS-F-P116-027}{\eta=0} and \NSi{NS-F-P116-028}{q=\tau}, so
\end{EESegment}
\NSFormula{NS-F-P116-029}{display}{none}
\[
u_{\theta,\mathrm{loc}}(\sqrt{2X_{\mathrm{in}}\tau},0,0,1-\tau)
=\tau^{-A}\bigl(e_0+O(\tau^{2h})\bigr).
\]
\begin{EESegment}{EE-TGT-P116-006}
This is Theorem~\ref{thm:3.1}\textup{(iv)}, and it completes the proof.
\end{EESegment}
\end{proof}

\begin{EESegment}{EE-TGT-P116-007}
\section{A compactly supported force and breakdown on all of space}\label{sec:10}
\end{EESegment}

\begin{EESegment}{EE-TGT-P116-008}
This section turns the fields from Theorem~\ref{thm:3.1} into the solution and force on all of
space in Theorem~\ref{thm:1.1}. Write the local fields as
\end{EESegment}
\NSFormula{NS-F-P116-030}{display}{none}
\[
u^{\mathrm{loc}}=\operatorname{curl}\mathcal A+\mathcal Be_\theta,
\qquad p^{\mathrm{loc}}.
\]
\begin{EESegment}{EE-TGT-P116-009}
Here \NSi{NS-F-P116-031}{\mathcal A} and \NSi{NS-F-P116-032}{\mathcal B} are the representatives
built in \eqref{eq:9.21}, and \NSi{NS-F-P116-033}{\mathcal B} does not depend on
\NSi{NS-F-P116-034}{\theta}. Start with viscosity one. The local velocity has the growth in
\eqref{eq:3.6} required by Theorem~\ref{thm:1.1}. What remains is to give it zero initial data and
compact spatial support while making its momentum residual agree, before time one, with a force in
\NSPage{117}%
\NSi{NS-F-P117-001}{C_c^\infty(\mathbb R^3\times(0,\infty);\mathbb R^3)}. This is done in
Proposition~\ref{prop:10.1} and Lemma~\ref{lem:10.3}.
\end{EESegment}

\begin{EESegment}{EE-TGT-P117-001}
Multiply the vector potential, the azimuthal coefficient, and the pressure by cutoffs that equal one
near \NSi{NS-F-P117-002}{(0,1)}. Lemma~\ref{lem:10.2} then proves that every derivative of the
new residual has a uniform limit as \NSi{NS-F-P117-003}{t\uparrow1}, and
Lemma~\ref{lem:10.3} builds a force for \NSi{NS-F-P117-004}{t>1} whose derivatives attain all
those limits. Before time one, the equation gives a uniform energy bound in
Lemma~\ref{lem:10.4}. Lemma~\ref{lem:10.5} then compares this field with any smooth solution of
bounded kinetic energy. That comparison transfers the growth of the local velocity to the smooth
solution and gives the contradiction needed for Theorem~\ref{thm:1.1}. Rescaling gives every
positive viscosity and also gives the periodic construction in Corollary~\ref{cor:10.6}.
\end{EESegment}

\begin{EESegment}{EE-TGT-P117-002}
\subsection{Keeping the construction inside the local domain.}\label{sec:10.1}
\end{EESegment}
\begin{EESegment}{EE-TGT-P117-003}
This subsection places the velocity and pressure inside the domain from
Theorem~\ref{thm:3.1}. First recall why the chosen vector potential
\NSi{NS-F-P117-005}{\mathcal A} is smooth at the axis, as stated in
Theorem~\ref{thm:3.1}\textup{(i)}, and why it vanishes in the heat exterior from
\eqref{eq:3.5}. Because of these two facts, its curl can be taken after the localization, and the
force at time one can be estimated in Lemma~\ref{lem:10.2}. For the background coefficients
\NSi{NS-F-P117-006}{U_n} from \eqref{eq:5.1}, the Stokes streamfunctions and their sum are
\end{EESegment}
\NSFormula{NS-F-P117-007}{display}{10.1}
\begin{equation}
\begin{gathered}
\mathcal S_n=q^{1-A+2nh}\int_0^XU_n(X',\eta)\,dX',
\qquad \mathcal S=\mathcal S_0+\sum_{n\geq1}\chi(c_nq)\mathcal S_n,\\
\mathcal A_{\mathrm{base}}=\frac{\mathcal S}r e_\theta.
\end{gathered}
\tag{10.1}\label{eq:10.1}
\end{equation}
\begin{EESegment}{EE-TGT-P117-004}
The cutoffs \NSi{NS-F-P117-008}{\chi(c_nq)} are the ones used in
Proposition~\ref{prop:5.5}. Equations~\eqref{eq:4.28} and \eqref{eq:5.10} give
\NSi{NS-F-P117-009}{\int_0^\infty U_n(X,\eta)\,dX=0} for every
\NSi{NS-F-P117-010}{n\geq0}. Each radial integral in \eqref{eq:10.1} is therefore zero beyond
the common support of the axial coefficients. Near the axis,
\NSi{NS-F-P117-011}{\mathcal S=r^2a(r^2,z,t)} for a smooth scalar
\NSi{NS-F-P117-012}{a}, so
\end{EESegment}
\NSFormula{NS-F-P117-013}{display}{none}
\[
(\mathcal S/r)e_\theta=a(r^2,z,t)(-x_2,x_1,0).
\]
\begin{EESegment}{EE-TGT-P117-005}
This is the smooth Cartesian form of the background potential.
\end{EESegment}

\begin{EESegment}{EE-TGT-P117-006}
Every wave velocity is the curl of its supported physical potential. For a prescribed axial mean
increment \NSi{NS-F-P117-014}{\gamma_{d,\mathrm{phys}}}, the mean construction uses the
azimuthal potential coefficient
\end{EESegment}
\NSFormula{NS-F-P117-015}{display}{none}
\[
\Psi=r^{-1}\mathcal I_c(r\gamma_{d,\mathrm{phys}}),
\qquad \mathcal I_cg=\mathcal Ig-\chi_m\mathcal Jg,
\]
\begin{EESegment}{EE-TGT-P117-007}
where the radial operations are defined in \eqref{eq:8.4}. Beyond both the source and the region
where \NSi{NS-F-P117-016}{\chi_m} changes, the identities
\NSi{NS-F-P117-017}{\chi_m=1} and \NSi{NS-F-P117-018}{\mathcal Ig=\mathcal Jg} hold, so
\NSi{NS-F-P117-019}{\Psi=0}. Together with their summation cutoffs, these wave and mean
potentials are the terms added to \NSi{NS-F-P117-020}{\mathcal A_{\mathrm{base}}} in
\eqref{eq:9.21}. The direct azimuthal mean increments are part of
\NSi{NS-F-P117-021}{\mathcal B}. Setting the auxiliary variables equal to
\NSi{NS-F-P117-022}{Y(r,t)} keeps the fields independent of
\NSi{NS-F-P117-023}{\theta}. Every annular correction vanishes near the axis for each fixed
\NSi{NS-F-P117-024}{t<1}. In particular, for the fixed
\NSi{NS-F-P117-025}{X_{\mathrm{ext}}} from Theorem~\ref{thm:3.1}, the chosen potential satisfies
\NSi{NS-F-P117-026}{\mathcal A=0} on \NSi{NS-F-P117-027}{X\geq X_{\mathrm{ext}}}.
\end{EESegment}

\begin{EESegment}{EE-TGT-P117-008}
Now choose cutoffs supported inside the local domain
\NSi{NS-F-P117-028}{\Omega_* = \{\tau>0,\ q<q_*\}}, where
\NSi{NS-F-P117-029}{\tau=1-t}. Apply the spatial cutoff before taking the curl, which keeps the
velocity incompressible. The time cutoff gives zero initial data.
\end{EESegment}

\begin{proposition}\label{prop:10.1}
\begin{EESegment}{EE-TGT-P117-009}
There are a compact set \NSi{NS-F-P117-030}{\mathcal K\subset\mathbb R^3}, a smooth vector field
\NSi{NS-F-P117-031}{u}, and a smooth scalar field \NSi{NS-F-P117-032}{p} on
\NSi{NS-F-P117-033}{\mathbb R^3\times[0,1)}. Their supports satisfy
\NSi{NS-F-P117-034}{\supp u(\mathord\cdot,t)\cup\supp p(\mathord\cdot,t)\subset\mathcal K}
for every \NSi{NS-F-P117-035}{0\leq t<1},
and they also satisfy
\end{EESegment}
\NSFormula{NS-F-P117-036}{display}{10.2}
\begin{equation}
\operatorname{div}u=0,
\qquad u(\mathord\cdot,t)=p(\mathord\cdot,t)=0
\qquad\text{for all sufficiently small }t\geq0.
\tag{10.2}\label{eq:10.2}
\end{equation}
\begin{EESegment}{EE-TGT-P117-010}
These fields agree with \NSi{NS-F-P117-037}{u^{\mathrm{loc}},p^{\mathrm{loc}}} on a spatial
neighborhood of the origin for every time sufficiently close to \NSi{NS-F-P117-038}{1}.
\end{EESegment}
\end{proposition}
\NSPage{118}%
\begin{proof}
\begin{EESegment}{EE-TGT-P118-001}
A fixed spatial cutoff will stay inside the local domain once \NSi{NS-F-P118-001}{q} is bounded
above in terms of \NSi{NS-F-P118-002}{z} and \NSi{NS-F-P118-003}{\tau}. The coordinates in
\eqref{eq:3.2} give such a bound uniformly in \NSi{NS-F-P118-004}{r}. If
\NSi{NS-F-P118-005}{1-\eta^2\geq1/2}, then \NSi{NS-F-P118-006}{q\leq2\tau}. Otherwise
\NSi{NS-F-P118-007}{|\eta|>2^{-1/2}}, and
\NSi{NS-F-P118-008}{q\leq(\sqrt2|z|)^{1/D}}. It follows that
\end{EESegment}
\NSFormula{NS-F-P118-009}{display}{10.3}
\begin{equation}
q\leq C_0\bigl(\tau+|z|^{1/D}\bigr).
\tag{10.3}\label{eq:10.3}
\end{equation}
\begin{EESegment}{EE-TGT-P118-002}
Choose \NSi{NS-F-P118-010}{0<\tau_0<1/2} and \NSi{NS-F-P118-011}{z_0>0} so that
\NSi{NS-F-P118-012}{C_0(\tau_0+z_0^{1/D})<q_*/2}. For any fixed
\NSi{NS-F-P118-013}{r_0>0}, choose a smooth axisymmetric cutoff
\NSi{NS-F-P118-014}{0\leq\chi_x\leq1} that is a smooth function of
\NSi{NS-F-P118-015}{r^2,z} and satisfies
\end{EESegment}
\NSFormula{NS-F-P118-016}{display}{none}
\[
\chi_x=1\quad(r\leq r_0/2,\ |z|\leq z_0/2),
\qquad \supp\chi_x\Subset\{r<r_0,\ |z|<z_0\}.
\]
\begin{EESegment}{EE-TGT-P118-003}
Also choose a smooth function \NSi{NS-F-P118-017}{0\leq\chi_t\leq1} of
\NSi{NS-F-P118-018}{t} with
\end{EESegment}
\NSFormula{NS-F-P118-019}{display}{none}
\[
\chi_t(t)=0\quad(1-t\geq\tau_0),
\qquad \chi_t(t)=1\quad(0\leq1-t\leq\tau_0/2),
\qquad c(x,t)=\chi_x(x)\chi_t(t).
\]
\begin{EESegment}{EE-TGT-P118-004}
By \eqref{eq:10.3}, the support of \NSi{NS-F-P118-020}{c} for
\NSi{NS-F-P118-021}{t<1} lies in \NSi{NS-F-P118-022}{q<q_*/2}. It therefore stays a fixed
distance, measured in the \NSi{NS-F-P118-023}{q} coordinate, from the outer boundary of
\NSi{NS-F-P118-024}{\Omega_*}. On that domain, define
\end{EESegment}
\NSFormula{NS-F-P118-025}{display}{10.4}
\begin{equation}
u=\operatorname{curl}(c\mathcal A)+c\mathcal Be_\theta,
\qquad p=cp^{\mathrm{loc}},
\tag{10.4}\label{eq:10.4}
\end{equation}
\begin{EESegment}{EE-TGT-P118-005}
and extend both fields by zero outside the support of the cutoff. The product rule gives
\end{EESegment}
\NSFormula{NS-F-P118-026}{display}{none}
\[
u=cu^{\mathrm{loc}}+\nabla c\times\mathcal A,
\qquad \operatorname{div}u=\operatorname{div}\operatorname{curl}(c\mathcal A)
 +r^{-1}\partial_\theta(c\mathcal B)=0.
\]
\begin{EESegment}{EE-TGT-P118-006}
The Cartesian forms recalled above make the fields smooth at the axis. The cutoffs, together with
the margin from \NSi{NS-F-P118-027}{q=q_*}, let the fields extend smoothly by zero everywhere
else. Their support lies in \NSi{NS-F-P118-028}{\mathcal K=\supp\chi_x}, and the time cutoff
gives \eqref{eq:10.2}. Both cutoffs equal one near
\NSi{NS-F-P118-029}{(0,1)}, so they leave the local fields unchanged there.
\end{EESegment}
\end{proof}

\begin{EESegment}{EE-TGT-P118-007}
For \NSi{NS-F-P118-030}{0\leq t<1}, define
\end{EESegment}
\NSFormula{NS-F-P118-031}{display}{10.5}
\begin{equation}
f=\mathscr R(u,p)=\partial_tu+(u\mathbin{\cdot}\nabla)u-\Delta u+\nabla p.
\tag{10.5}\label{eq:10.5}
\end{equation}
\begin{EESegment}{EE-TGT-P118-008}
This definition includes every derivative of the cutoffs. Taking its divergence gives
\end{EESegment}
\NSFormula{NS-F-P118-032}{display}{none}
\[
-\Delta p=\sum_{i,j=1}^3\partial_i\partial_j(u_iu_j)-\operatorname{div}f.
\]
\begin{EESegment}{EE-TGT-P118-009}
The localized pressure therefore satisfies the Poisson equation with the matching term from the
divergence of the force. The force is allowed to have nonzero divergence.
\end{EESegment}

\begin{EESegment}{EE-TGT-P118-010}
\subsection{Limits of force derivatives at the final time.}\label{sec:10.2}
\end{EESegment}
\begin{EESegment}{EE-TGT-P118-011}
So far, the force in \eqref{eq:10.5} is defined only for
\NSi{NS-F-P118-033}{t<1}. To extend it smoothly past time one, all of its derivatives must have
compatible limits as \NSi{NS-F-P118-034}{t\uparrow1}. At the origin, those limits are controlled
by the flatness of the local residual in \eqref{eq:3.4}. Away from the origin, the terms introduced
by localization are controlled by the endpoint bounds in Theorem~\ref{thm:3.1}\textup{(ii)} and
the exact exterior heat flow in \eqref{eq:3.5}. Lemma~\ref{lem:10.2} first proves the derivative
limits. Lemma~\ref{lem:10.3} then builds an extension with compact support in time.
\end{EESegment}

\begin{lemma}\label{lem:10.2}
\begin{EESegment}{EE-TGT-P118-012}
For every spatial multi-index \NSi{NS-F-P118-035}{\alpha} and every integer
\NSi{NS-F-P118-036}{j\geq0}, the derivative
\NSi{NS-F-P118-037}{\partial_x^\alpha\partial_t^jf} converges uniformly on
\NSi{NS-F-P118-038}{\mathbb R^3} as \NSi{NS-F-P118-039}{t\uparrow1}. There are functions
\NSi{NS-F-P118-040}{F_j\in C_c^\infty(\mathbb R^3;\mathbb R^3)}, all supported in
\NSi{NS-F-P118-041}{\mathcal K}, such that
\end{EESegment}
\NSFormula{NS-F-P118-042}{display}{10.6}
\begin{equation}
\lim_{t\uparrow1}\partial_x^\alpha\partial_t^jf(x,t)=\partial_x^\alpha F_j(x),
\qquad \partial_x^\alpha F_j(0)=0.
\tag{10.6}\label{eq:10.6}
\end{equation}
\begin{EESegment}{EE-TGT-P118-013}
These are the derivative limits of \NSi{NS-F-P118-043}{f} from
\NSi{NS-F-P118-044}{t<1}, and they agree with both spatial and time differentiation.
\end{EESegment}
\end{lemma}
\NSPage{119}%
\begin{proof}
\begin{EESegment}{EE-TGT-P119-001}
Start with endpoint bounds away from the origin. Take a compact spatial set on which
\NSi{NS-F-P119-001}{q} has a positive lower bound and stays away from
\NSi{NS-F-P119-002}{q_*}. Theorem~\ref{thm:3.1}\textup{(ii)} bounds every physical derivative
of the actual potentials and scalar fields there. More exactly, let
\NSi{NS-F-P119-003}{K_0} be such a compact set, let
\NSi{NS-F-P119-004}{G=\mathcal A,\mathcal Be_\theta,p^{\mathrm{loc}}}, and take
\NSi{NS-F-P119-005}{t<t'<1}. The fundamental theorem of calculus gives
\end{EESegment}
\NSFormula{NS-F-P119-006}{display}{none}
\[
\sup_{K_0}|\partial_x^\alpha\partial_t^jG(\mathord\cdot,t')
 -\partial_x^\alpha\partial_t^jG(\mathord\cdot,t)|
\leq|t'-t|\sup_{K_0\times[t,t']}|\partial_x^\alpha\partial_t^{j+1}G|.
\]
\begin{EESegment}{EE-TGT-P119-002}
The last supremum has a bound near time one that does not depend on
\NSi{NS-F-P119-007}{t,t'}. Every fixed derivative is therefore uniformly Cauchy as
\NSi{NS-F-P119-008}{t\uparrow1}. These constants may depend on the chosen positive lower bound
for \NSi{NS-F-P119-009}{q}.
\end{EESegment}

\begin{EESegment}{EE-TGT-P119-003}
The argument in Proposition~\ref{prop:9.9} that leaves only finitely many nonzero cutoff terms
gives these bounds on the closed parameter range
\NSi{NS-F-P119-010}{-1\leq\eta\leq1}. It gives them for the actual representatives and for the
smooth Cartesian factors in their near-axis forms, while the labels and integer frequencies stay
fixed. Multiplying by the fixed cutoffs and applying the finite differential expression
\eqref{eq:10.5} keep the bounds.
\end{EESegment}

\begin{EESegment}{EE-TGT-P119-004}
In the region that remains, \NSi{NS-F-P119-011}{r} has a positive lower bound and
\NSi{NS-F-P119-012}{q} is small enough that
\NSi{NS-F-P119-013}{X\geq X_{\mathrm{ext}}}. The velocity there is the exterior azimuthal field
\NSi{NS-F-P119-014}{K(r,\tau)e_\theta} from \eqref{eq:3.5}. In terms of the heat profile
\NSi{NS-F-P119-015}{\mathcal H} defined in \eqref{eq:A.32},
\end{EESegment}
\NSFormula{NS-F-P119-016}{display}{10.7}
\begin{equation}
K(r,\tau)=c_\infty s^{-A}\mathcal H(2\tau/s),
\qquad s=r^2/2.
\tag{10.7}\label{eq:10.7}
\end{equation}
\begin{EESegment}{EE-TGT-P119-005}
Thus the function in Theorem~\ref{thm:3.1} is
\NSi{NS-F-P119-017}{\mathcal H_{\mathrm{ext}}(Z)=2^Ac_\infty\mathcal H(4Z)}. For each
\NSi{NS-F-P119-018}{j}, the integral in \eqref{eq:A.32} can be differentiated uniformly for
\NSi{NS-F-P119-019}{Z\geq0} because its differentiated integrand is bounded by a constant times
\NSi{NS-F-P119-020}{e^{-v}v^{h+j}}. It follows that
\end{EESegment}
\NSFormula{NS-F-P119-021}{display}{10.8}
\begin{equation}
|\partial_\tau^jK(r,\tau)|\leq C_jr^{-1-2h-2j},
\qquad
\left|\partial_\tau^j\frac{K(\rho,\tau)^2}{\rho}\right|
\leq C_j\rho^{-3-4h-2j}.
\tag{10.8}\label{eq:10.8}
\end{equation}
\begin{EESegment}{EE-TGT-P119-006}
The pressure is normalized to vanish at radial infinity. On a compact interval of positive radii
\NSi{NS-F-P119-022}{r\in[r_-,r_+]}, where \NSi{NS-F-P119-023}{r_->0}, the second upper bound
is integrable over \NSi{NS-F-P119-024}{\rho\geq r_-}. Dominated convergence therefore gives
\end{EESegment}
\NSFormula{NS-F-P119-025}{display}{none}
\[
\partial_\tau^jp^{\mathrm{loc}}(r,\tau)
=-\int_r^\infty\partial_\tau^j\left(\frac{K(\rho,\tau)^2}{\rho}\right)\,d\rho.
\]
\begin{EESegment}{EE-TGT-P119-007}
It also gives a uniform bound and a one-sided limit for every derivative on each compact interval of
positive radii. Here \NSi{NS-F-P119-026}{\partial_t=-\partial_\tau}. Differentiate the lower
endpoint to get the higher radial derivatives of the pressure, and differentiate \eqref{eq:10.7}
directly to get those of \NSi{NS-F-P119-027}{K}. This gives the required smooth one-sided
extension to \NSi{NS-F-P119-028}{\tau=0}, using the integral on
\NSi{NS-F-P119-029}{\tau\geq0}. In this region \NSi{NS-F-P119-030}{\mathcal A=0}, so after
localization the fields contain only these smooth exterior quantities and the fixed cutoffs. The
heat equation for \NSi{NS-F-P119-031}{K} and the identity
\NSi{NS-F-P119-032}{\partial_rp^{\mathrm{loc}}=K^2/r} also show directly that the residual of the
uncut exterior field is exactly zero.
\end{EESegment}

\begin{EESegment}{EE-TGT-P119-008}
These two regions cover the final time slice away from the origin. If
\NSi{NS-F-P119-033}{z\neq0}, then \NSi{NS-F-P119-034}{q\geq|z|^{1/D}>0}. If
\NSi{NS-F-P119-035}{z=0} and \NSi{NS-F-P119-036}{x\neq0}, then
\NSi{NS-F-P119-037}{r>0}, and the exterior argument above applies on a small neighborhood.
Every Cartesian space--time derivative of the force therefore converges uniformly on compact
spatial sets that do not contain the origin.
\end{EESegment}

\begin{EESegment}{EE-TGT-P119-009}
Near the origin, both cutoffs equal one. Use \eqref{eq:3.4} on
\NSi{NS-F-P119-038}{X\leq X_{\mathrm{ext}}}, and use the exactly zero residual when
\NSi{NS-F-P119-039}{X} is larger. For every
\NSi{NS-F-P119-040}{\alpha,j,N}, this gives, on one fixed neighborhood of
\NSi{NS-F-P119-041}{(0,1)},
\end{EESegment}
\NSFormula{NS-F-P119-042}{display}{10.9}
\begin{equation}
|\partial_x^\alpha\partial_t^jf(x,t)|
\leq C_{\alpha,j,N}\bigl(\tau+|z|^{1/D}\bigr)^N.
\tag{10.9}\label{eq:10.9}
\end{equation}
\begin{EESegment}{EE-TGT-P119-010}
Given any tolerance, first choose a small ball and then choose a time close enough to one that this
bound is below the tolerance throughout the ball. A finite cover of the part of
\NSi{NS-F-P119-043}{\mathcal K} outside that ball gives uniform convergence there. Every one of
these derivatives is zero outside \NSi{NS-F-P119-044}{\mathcal K}. They are therefore uniformly
Cauchy on \NSi{NS-F-P119-045}{\mathbb R^3}, and their limits are zero at the origin.
\end{EESegment}
\NSPage{120}%
\begin{EESegment}{EE-TGT-P120-001}
Define \NSi{NS-F-P120-001}{F_j} as the uniform limit for
\NSi{NS-F-P120-002}{\alpha=0}. Pass to the limit in the fundamental theorem of calculus along
segments parallel to the spatial coordinate axes. This identifies every other limit as
\NSi{NS-F-P120-003}{\partial_x^\alpha F_j}. It follows that
\NSi{NS-F-P120-004}{F_j} is smooth and has the stated support. Passing to the limit in the
matching identity in time gives
\end{EESegment}
\NSFormula{NS-F-P120-005}{display}{10.10}
\begin{equation}
\partial_x^\alpha\partial_t^jf(x,t)
=\partial_x^\alpha F_j(x)-\int_t^1\partial_x^\alpha\partial_t^{j+1}f(x,s)\,ds.
\tag{10.10}\label{eq:10.10}
\end{equation}
\begin{EESegment}{EE-TGT-P120-002}
This proves that the derivative limits agree with both spatial and time differentiation, and it
finishes the proof of \eqref{eq:10.6}.
\end{EESegment}
\end{proof}

\begin{EESegment}{EE-TGT-P120-003}
The limits \NSi{NS-F-P120-006}{F_j} are the Taylor coefficients needed to continue the force
through time one. Give the coefficients time cutoffs whose supports shrink as the derivative order
grows, so the extension also has compact support in time. A smooth force with compact support has
the decay required in \cite[(5)]{ref-13}, and the zero initial data has the required initial-data
properties as well.
\end{EESegment}

\begin{lemma}\label{lem:10.3}
\begin{EESegment}{EE-TGT-P120-004}
The force in \eqref{eq:10.5} extends to
\NSi{NS-F-P120-007}{f\in C_c^\infty(\mathbb R^3\times(0,\infty);\mathbb R^3)}, with support in
\NSi{NS-F-P120-008}{\mathcal K\times[0,2]}.
\end{EESegment}
\end{lemma}

\begin{proof}
\begin{EESegment}{EE-TGT-P120-005}
Choose \NSi{NS-F-P120-009}{\chi_0\in C_c^\infty(\mathbb R)} so that it equals one near zero and
equals zero whenever its argument is at least one. For
\NSi{NS-F-P120-010}{\sigma=t-1\geq0}, set
\end{EESegment}
\NSFormula{NS-F-P120-011}{display}{10.11}
\begin{equation}
f(x,1+\sigma)=\sum_{j=0}^\infty\chi_0(b_j\sigma)\frac{\sigma^j}{j!}F_j(x),
\tag{10.11}\label{eq:10.11}
\end{equation}
\begin{EESegment}{EE-TGT-P120-006}
where the increasing integers \NSi{NS-F-P120-012}{b_j\geq1} are chosen below. On the support of
each summand, \NSi{NS-F-P120-013}{0\leq\sigma\leq b_j^{-1}}. The product rule then gives
\end{EESegment}
\NSFormula{NS-F-P120-014}{display}{10.12}
\begin{equation}
\left\|\partial_x^\alpha\partial_\sigma^m
 \left(\chi_0(b_j\sigma)\frac{\sigma^j}{j!}F_j\right)\right\|_\infty
\leq C_{j,m}b_j^{m-j}\|F_j\|_{C^{|\alpha|}}.
\tag{10.12}\label{eq:10.12}
\end{equation}
\begin{EESegment}{EE-TGT-P120-007}
To see this, suppose \NSi{NS-F-P120-015}{a} derivatives land on the cutoff. Their factor
\NSi{NS-F-P120-016}{b_j^a} is multiplied by a monomial of degree
\NSi{NS-F-P120-017}{j-m+a}, which is bounded by
\NSi{NS-F-P120-018}{b_j^{-(j-m+a)}}. If that monomial has been differentiated to zero, the term
is absent. Set \NSi{NS-F-P120-019}{b_0=1}. For \NSi{NS-F-P120-020}{j\geq1}, define
\end{EESegment}
\NSFormula{NS-F-P120-021}{display}{none}
\[
b_j=\min\left\{b\in\mathbb N:b>b_{j-1},\quad
\max_{\substack{a,m\geq0\\a+m\leq\lfloor j/2\rfloor}}
C_{j,m}b^{m-j}\|F_j\|_{C^a}\leq2^{-j}\right\},
\]
\begin{EESegment}{EE-TGT-P120-008}
where the maximum is taken over nonnegative integers \NSi{NS-F-P120-022}{a,m}. This gives only
finitely many conditions, and each one has \NSi{NS-F-P120-023}{m-j<0}. So the defining set
contains at least one integer. With this choice, \eqref{eq:10.12} is at most
\NSi{NS-F-P120-024}{2^{-j}} whenever
\NSi{NS-F-P120-025}{|\alpha|+m\leq\lfloor j/2\rfloor}.
\end{EESegment}

\begin{EESegment}{EE-TGT-P120-009}
After any fixed number of derivatives, the tail now converges uniformly. At
\NSi{NS-F-P120-026}{\sigma=0}, the cutoff in each fixed summand is constant on a neighborhood of
zero. Its \NSi{NS-F-P120-027}{m}-th derivative therefore contributes only when
\NSi{NS-F-P120-028}{j=m}, and that contribution is \NSi{NS-F-P120-029}{F_m}. Uniform
convergence after differentiation, together with \eqref{eq:10.10}, shows that \eqref{eq:10.11}
matches every mixed space--time derivative limit from \NSi{NS-F-P120-030}{t<1}. Every summand
keeps its support in \NSi{NS-F-P120-031}{\mathcal K} and is zero when
\NSi{NS-F-P120-032}{\sigma\geq1}. The same uniform convergence proves smoothness at those support
boundaries. Finally, the original force is zero near
\NSi{NS-F-P120-033}{t=0}, by \eqref{eq:10.2} and the construction of
\NSi{NS-F-P120-034}{p}.
\end{EESegment}

\begin{EESegment}{EE-TGT-P120-010}
For the decay bounds, choose \NSi{NS-F-P120-035}{R} with
\NSi{NS-F-P120-036}{\mathcal K\subset B(0,R)}, and write
\NSi{NS-F-P120-037}{M_{\alpha,m}=\|\partial_x^\alpha\partial_t^mf\|_\infty}. Since the force
has compact support, every integer \NSi{NS-F-P120-038}{k\geq0} satisfies
\end{EESegment}
\NSFormula{NS-F-P120-039}{display}{none}
\[
|\partial_x^\alpha\partial_t^mf(x,t)|
\leq M_{\alpha,m}(3+R)^k(1+|x|+t)^{-k}.
\]
\end{proof}
